\begin{abstract}
%% Camera-ready rewrite (comparative-study framing). Submitted text: sec/0_abstract.tex.
\begin{revpar}
Vision foundation models such as CLIP and the Segment Anything Model (SAM) are
increasingly coupled through a vision--language correlation map for
open-vocabulary segmentation, yet it is unclear which parts of such pipelines
carry over to class-agnostic, binary salient object detection (SOD). We present
\methodname{}, an adaptation of the SAM--CLIP correlation paradigm to SOD built on
a CLIP ViT-L/14 encoder, a bank of salient-object prompts, and a pseudo-prompt
generator with an adaptive fallback that never returns an empty prompt set.
Around this model we conduct a controlled comparative study of the adaptation
choices: prompt bank size, wording and aggregation, a prompt-residual connection,
multi-token text fusion, decoder and loss design, and encoder fine-tuning,
evaluated on six benchmarks with repeated runs. Fine-tuning the CLIP visual encoder
is the dominant factor, improving S-measure by $0.8$--$4.4$ points over frozen
encoders, whereas the prompt-side components add at most $0.3$ points and several
popular choices have no consistent effect. With RGB input only, \methodname{}
obtains the best S-measure among the compared methods on four of five RGB
benchmarks and outperforms recent RGB-D methods on SIP. Code is available at
\url{https://github.com/aminullah6264/CLIPSAM}.
\end{revpar}
\end{abstract}
